% =========================================================
% Proof Stub: Absolute Convergence Implies Convergence
% Source: volume-iii/analysis/sequences/notes/series/notes-absolute-convergence.tex
% Mode: standard
% =========================================================

\clearpage
\phantomsection
\hypertarget{proof-RA-ABB-P-abs-conv}{}

\subsubsection[Absolute Convergence Implies Convergence]{Proof --- Absolute Convergence Implies Convergence}

\bigskip

\noindent
\textbf{Source.}
Abbott, \emph{Understanding Analysis}, Theorem~2.7.6.
See notes \texttt{notes-absolute-convergence.tex}.

\vspace{0.75em}

\noindent
\textbf{Goal.}
If $\sum_{n=1}^\infty |a_n|$ converges, then $\sum_{n=1}^\infty a_n$
converges.

\vspace{0.75em}

\noindent
\textbf{Logical form.}
\[
  \sum |a_n| < \infty
  \;\Rightarrow\;
  \sum a_n \text{ converges.}
\]

\vspace{0.75em}

\noindent
\textbf{Proof strategy.}
Let $S_N = \sum_{n=1}^N a_n$. Show $(S_N)$ is Cauchy.
For $m > n$: $|S_m - S_n| = |\sum_{k=n+1}^m a_k| \le \sum_{k=n+1}^m |a_k|$.
Since $\sum |a_k|$ converges, its partial sums are Cauchy, so the tail
$\sum_{k=n+1}^m |a_k| \to 0$. By Cauchy Criterion, $(S_N)$ converges.

\vspace{0.75em}

\noindent
\textbf{Proof.}
\begin{proof}
\Claim $\sum |a_n| < \infty \Rightarrow \sum a_n$ converges.

\Given $\sum_{n=1}^\infty |a_n|$ converges.

\Goal Show partial sums $S_N = \sum_{n=1}^N a_n$ are Cauchy.

\Strategy Bound $|S_m - S_n|$ by a tail of the absolutely convergent series.

\medskip

% --- let T_N = Σ_{n=1}^N |a_n|; (T_N) is Cauchy (convergent) ---

% --- for m > n: |S_m - S_n| ≤ Σ_{k=n+1}^m |a_k| = T_m - T_n ---

% --- (T_N) Cauchy: for ε > 0, ∃ N with T_m - T_n < ε for m > n ≥ N ---

% --- hence |S_m - S_n| < ε; (S_N) Cauchy; apply Cauchy Criterion ---

\end{proof}

\begin{remark*}[Proof shape]
Triangle inequality on partial sums, then Cauchy Criterion for sequences.
\end{remark*}

\begin{remark*}[Dependencies]
\begin{itemize}
  \item Cauchy Criterion in $\mathbb{R}$.
  \item Triangle inequality for finite sums.
  \item Convergent sequence implies Cauchy (applied to $(T_N)$).
\end{itemize}
\end{remark*}
